\section{Introduction}
\label{sec:intro}

A PDF file is a program for a printer: a page tree, content streams, fonts, and a
cross-reference index that says where every object lives. Readers are therefore easy to
classify by what they do with those bytes. A \emph{renderer} reads them. A \emph{reader with
comments} either writes new objects into the file --- which changes the bytes other
applications see --- or keeps annotations in a companion file, which the reader next to it
will never show.

The distinction matters because a PDF is usually not the only copy. An annotated manuscript
is mailed to a co-author, a marked-up drawing goes to a fabricator, a reviewed contract goes
back to a lawyer. An annotation that lives in a sidecar reaches none of them.

Writing into a PDF is not hard to begin and hard to get right. The format is designed for
\emph{incremental} updates: append objects, rewrite the dictionaries that changed, add
another cross-reference section pointing back at the previous one, and the file remains valid
for every existing reader. The difficulty is that an update is a \emph{diff}, so every save
has to know what the previous save did. A writer that only appends is correct exactly once.

This paper reports what it took to build such a reader as a self-contained system:
\textbf{PDF Looker Pro}, an MIT-licensed desktop application for Windows whose PDF layer
--- parser, incremental writer, whole-document rewrite and merger --- is written from
scratch in TypeScript. The only runtime dependency is pdf.js~\cite{pdfjs}, used for
rasterisation and text extraction; the page model, the writers and the annotation model are
the authors' own. The feature set follows that of a long-standing commercial Windows viewer
for which development has ceased~\cite{pdfxchange}; there is no affiliation with it and no
code in common.

\paragraph{Contributions.} We claim no new algorithm. The contributions are an account of
the invariants that a system of this kind must satisfy, the design that satisfies them, and a
verification methodology with an honest record of what it caught and what it missed.

\begin{enumerate}
  \item \textbf{A deletion-correct incremental annotation writer}
        (\S\ref{sec:incremental}). A page's \texttt{/Annots} is rebuilt as
        \emph{(existing $\setminus$ replaced) $\cup$ new}, where ``replaced'' is the set of
        object numbers the previous save produced. A page whose annotations were all deleted
        is still rewritten, because nothing else removes an obsolete entry. We show the two
        user-visible symptoms of violating this invariant --- annotations duplicating on every
        save, and deleted annotations reappearing --- both of which occurred in this codebase
        and both of which are now covered by a regression test.

  \item \textbf{A whole-document rewrite that materialises inheritance}
        (\S\ref{sec:rewrite}). Flattening the page tree to support delete, insert, reorder,
        extract and crop destroys attributes a page inherited from intermediate
        \texttt{/Pages} nodes (\texttt{MediaBox}, \texttt{CropBox}, \texttt{Resources},
        \texttt{Rotate}, \texttt{UserUnit}). The rewrite writes them onto every page, copies
        objects verbatim when they are stored uncompressed, re-serialises those inside object
        streams, and recomputes the cross-reference.

  \item \textbf{Cross-document merging by reference renumbering}
        (\S\ref{sec:merge}). The incoming object graph is renumbered above the host's highest
        object number; every reference is rewritten, every stream's bytes are copied
        unchanged, and the incoming pages are rebuilt on the host's page tree.

  \item \textbf{A dual-representation annotation model} (\S\ref{sec:annot-model}) that writes
        a standard subtype so other readers display the mark, and a private dictionary so this
        reader restores it without loss. We state the cost: two representations can disagree,
        and the standard one cannot express everything the private one can.

  \item \textbf{Search over a virtualised document} (\S\ref{sec:search}) that highlights by
        rewriting the text layer's spans, so a painted match is aligned with the painted
        glyphs by construction, guarded by a comparison against the strings the index was
        built from.

  \item \textbf{A dual-target verification methodology} (\S\ref{sec:verification}): 789
        behavioural assertions over 31 suites, 734 driving a browser build and 55 driving the
        packaged Electron application as a real executable, with a case study of a defect the
        second target caught and of three defects the authors introduced into the harness
        itself.
\end{enumerate}

\paragraph{Scope.} The paper describes a desktop application, not a library or a format
extension. It does not claim conformance to any profile of ISO 32000~\cite{iso32000}, does not
handle encrypted documents for writing, and does not implement forms, digital signatures or
optical character recognition. Section~\ref{sec:limits} lists the limitations in full; we
regard that list as a deliverable, not a disclaimer.

\paragraph{Availability.} Source, tests, build scripts and the measurement scripts that
produced every number in this paper are at \repo, MIT licensed. The packaged Windows build is
published as a release artefact of the same repository.
